% 360_HRV_Galois_Experiment_En_ch.tex
% Johann Pascher, 8 September 2026
% Galois-Informed HRV Analysis: Implementation and Exploratory Assessment
\providecommand{\BM}{\textbf{[B]}}
\providecommand{\KM}{\textbf{[K]}}
\providecommand{\SM}{\textbf{[S]}}
\providecommand{\XM}{\textbf{[X]}}
\providecommand{\LM}{\textbf{[L]}}

\phantomsection
\addcontentsline{toc}{chapter}{%
  Galois-Informed HRV Analysis --- Implementation and Exploratory Assessment}

\section*{Status markers}
\begin{center}
\begin{tabular}{@{}lp{10.5cm}@{}}
\textbf{[B]} & \emph{Proved} --- algebraically proved. \\[3pt]
\textbf{[K]} & \emph{Confirmed} --- numerically verified. \\[3pt]
\textbf{[S]} & \emph{Speculative} --- open. \\[3pt]
\textbf{[L]} & \emph{Literature evidence} --- external primary source. \\
\end{tabular}
\end{center}

\section*{Abstract}

Docs.~358 and 359 formulate an algebraic hypothesis for HRV frequency
analysis: frequency ratios of coupled physiological oscillators should lie
on a Galois grid, with the structural tolerance $100\xi\approx1.33\,\%$.
This document implements the approach as a Python module and assesses it
exploratively on two public ECG datasets and on the author's own Polar H10
recordings. The result is negative: the present data do not support the
hypothesis, and they cannot test it either, because LF/HF is a power
ratio, not a frequency ratio.

\textbf{Results:}
\begin{enumerate}
\item \textbf{Module} \KM: \texttt{galois\_hrv\_preprocess.py} ---
  Galois projection, HRV band limits, orbit depth, complexity measure
  $2/(p+q)$.
\item \textbf{Null model} \KM: In the grid up to denominator 60, every
  value between 1 and 4 lies within 1.33\,\% of a grid element; 39\,\% of
  random values lie within 0.06\,\%. Close grid hits are therefore no
  evidence on their own.
\item \textbf{BIDMC-01} \LM: intensive-care patient, 88~y, 8~min. No
  ratio meets the 5-limit criterion; the dataset is unsuitable for the
  question.
\item \textbf{Fantasia} \LM: 5 young and 5 old healthy subjects, about
  120~min each, spontaneous breathing. LF peak CV 11--33\,\%; 5-limit
  LF/HF 4/10 against a null rate of 21\,\% (binomial test $p=0.14$).
\item \textbf{Polar H10} \KM: With validated beat detection, the dominant
  spectral peak during paced breathing (4--6/min) lies within one frequency
  bin of the breathing frequency, and the HF peak is its second harmonic.
  Grid hits of LF/HF occur only at chance frequency.
\item \textbf{Structural reason} \BM: Under spontaneous breathing the HF
  peak is the respiratory frequency and LF/HF the ratio of two independent
  control loops; with paced breathing below 9/min, respiration itself lies
  in the LF band.
\item \textbf{Next step} \SM: prospective test of frequency coupling
  during paced breathing (Section~\ref{sec:360next}).
\end{enumerate}

% ============================================================
\section{Galois HRV preprocessing module}
% ============================================================

\texttt{galois\_hrv\_preprocess.py} implements (Doc.~359 §5):

\begin{description}
\item[Galois grid] 901 rational ratios $p/q\le4$ with prime factors
  $\subseteq\{2,3,5,7,11,13\}$ up to denominator 60. The prime set is
  $\{3\}\cup\{p : p\mid 3^k-1,\ k\le6\}$; a prime $p\neq3$ belongs to
  depth $k$ when $k$ is a multiple of $\mathrm{ord}_p(3)$.
\item[Tolerance] $100\xi = 1/75 \approx 1.33\,\%$ with Farey height
  $Q_c=\pi/\sqrt{6\cdot100\xi}\approx11.1$ (Doc.~343): at this denominator
  the mean Farey spacing equals twice the tolerance. A grid up to
  denominator 60 is denser than the tolerance and does not separate.
\item[HRV band limits] VLF: $[0,\,\tfrac{1}{25}]$~Hz, LF: $[\tfrac{1}{25},\,\tfrac{3}{20}]$~Hz,
  HF: $[\tfrac{3}{20},\,\tfrac{2}{5}]$~Hz. The prime factors lie in
  $\{2,3,5\}$ \BM; the limits are thus compatible with the 5-limit grid,
  but not determined by it (Section~\ref{sec:360baender}).
\item[Galois projection] ratio $\mapsto$ nearest $p/q$ in the grid,
  distance $d=|r/(p/q)-1|$.
\item[Orbit depth] $k=\max\,\mathrm{ord}_p(3)$ over the prime factors
  $p\notin\{2,3\}$ (Docs.~343, 358).
\item[Complexity measure] $2/(p+q)$ orders ratios by their simplicity
  (Doc.~328). It is not a computed Arnold tongue width; that would require
  the coupling strength of a specific oscillator model.
\item[Pipeline] \texttt{hrv\_galois\_pipeline(rr\_ms)} --- one function call.
\end{description}

\textbf{Resolution condition.} A classification with the tolerance
$\delta=100\xi$ is meaningful only if the frequency resolution satisfies
$\Delta f/f\le\delta$. With $\Delta f=1/T_w$ for the window length $T_w$
(in Welch's method the segment length, not the recording duration) this
gives $T_w\ge1/(\delta f)$: 12.5~min at 0.10~Hz, 15~min at 0.083~Hz \BM.
This is distinct from the minimum separation of two resolvable peaks
(about $2/T_w$ for a Hann window) and from the precision of the position
of a single peak, which interpolation can make finer than $\Delta f$.

\textbf{Frequency and power ratios.} The hypothesis concerns frequency
ratios of coupled oscillators. LF/HF is a ratio of integrated band powers;
there is no reason to expect it on the grid. In what follows, the
projection of LF/HF serves only to show the behaviour of the null model.

% ============================================================
\section{Band limits and null model}\label{sec:360baender}
% ============================================================

\textbf{Band limits} \KM. The three Task Force limits have prime factors in
$\{2,3,5\}$ in Hz, mHz, cycles per minute ($2.4=12/5$, 9, 24) and as
periods (25~s, $20/3$~s, $5/2$~s), because the conversion factors 60 and
1000 contain only these primes. In rad/s they are irrational and cannot be
classified. The property is not specific either: 8 of the 9 numbers with
one significant digit and a third of all two-decimal values in $(0,1)$ are
5-limit. Rounded conventions therefore land on the grid by themselves; the
band limits are compatible, but not algebraically determined. Frequency
ratios, by contrast, are dimensionless and independent of units.

\textbf{Null model} \KM. Probability that a log-uniformly distributed
value in $[1,4]$ lies within the tolerance $\tau$ of a grid element
(200\,000 draws):

\begin{center}
\begin{tabular}{@{}lrr@{}}
\toprule
$\tau$ & Grid up to denominator 60 (901) & Grid up to denominator 11 (133) \\
\midrule
1.33\,\% & 100.0\,\% & 90.5\,\% \\
0.14\,\% & 73.9\,\% & 18.5\,\% \\
0.06\,\% & 39.5\,\% & 7.9\,\% \\
0.02\,\% & 14.5\,\% & 2.6\,\% \\
\bottomrule
\end{tabular}
\end{center}

With the grid up to denominator 60, the criterion ``within 1.33\,\%''
cannot distinguish anything. Even the grid up to the Farey height 11 still
covers 91\,\% at 1.33\,\%.

% ============================================================
\section{Negative result I: BIDMC-01}
% ============================================================

\textbf{Dataset:} BIDMC PPG and Respiration Dataset v1.0.0,
record 01 \LM\,[Pim18].
88~y, male, MICU, 8~min, 125~Hz PPG, ECG and respiration.

\textbf{Result:} 0/3 frequency ratios meet the 5-limit criterion (Welch
and Lomb-Scargle). VLF cannot be resolved in 8~min. DFA $\alpha_1=0.50$,
$\alpha_2=0.45$ --- white noise; in healthy subjects $\approx1.0$.

\textbf{Why:}
\begin{enumerate}
\item Too short for VLF.
\item Intensive-care patient: fractal regulation largely lost.
\item HF peak = measured respiratory frequency (0.328~Hz versus 0.336~Hz in
  the RESP channel).
\item Measurement uncertainty Welch versus Lomb-Scargle $\pm2\,\%$, larger
  than the tolerance of 1.33\,\%.
\end{enumerate}

% ============================================================
\section{Negative result II: Fantasia}
% ============================================================

\textbf{Dataset:} PhysioNet Fantasia Database v1.0.0
\LM\,[Iye96]. 5 young (f1y01--05, 21--34~y) + 5 old
(f2o01--05, 68--85~y), about 120~min rest each, 250~Hz ECG,
spontaneous breathing. Reference beat annotations (.ecg), CC-BY 4.0.

\textbf{Method:} wfdb annotations, Welch 10-min windows,
zero-padding $\times4$, parabolic peak interpolation,
LF stability over 6 $\times$ 5-min segments.

\begin{center}
\begin{tabular}{@{}lrrrrrrrl@{}}
\toprule
Rec & $n_\text{RR}$ & LF~Hz & CV\,\% & LF/HF & $p/q$ & $k$ & 5-lim & $2/(p+q)$ \\
\midrule
f1y01 & 8707 & 0.109 & 11.5 & 0.438 & $7/16$ & 6 & N & 0.087 \\
f1y02 & 6919 & 0.085 & 33.0 & 0.867 & $13/15$ & 4 & N & 0.071 \\
f1y03 & 7641 & 0.100 & 13.7 & 4.219 & $4$ & 1 & Y & 0.400 \\
f1y04 & 5232 & 0.078 & 30.9 & 1.060 & $35/33$ & 6 & N & 0.029 \\
f1y05 & 6943 & 0.068 & 16.5 & 2.993 & $3$ & 1 & Y & 0.500 \\
\midrule
f2o01 & 7045 & 0.067 & 25.6 & 1.494 & $3/2$ & 1 & Y & 0.400 \\
f2o02 & 6327 & 0.052 & 14.4 & 0.864 & $45/52$ & 4 & N & 0.021 \\
f2o03 & 6478 & 0.055 & 22.0 & 0.841 & $21/25$ & 6 & N & 0.043 \\
f2o04 & 6849 & 0.065 & 32.2 & 0.719 & $28/39$ & 6 & N & 0.030 \\
f2o05 & 8334 & 0.049 & 11.4 & 1.873 & $15/8$ & 4 & Y & 0.087 \\
\bottomrule
\end{tabular}
\end{center}

\textbf{Statistical tests (Mann-Whitney~U):} All group differences young
versus old n.s.\ ($p>0.40$).

\textbf{Findings:}
\begin{itemize}
\item LF peak stability: 0/10 records stable ($\text{CV}>1.33\,\%$).
\item 5-limit LF/HF (nearest grid element has prime factors in
  $\{2,3,5\}$): 4/10. The null rate for this criterion is 21\,\%;
  one-sided binomial test $p=0.14$, i.e.\ not above chance.
\item No group difference young/old ($n=5$ per group).
\end{itemize}

% ============================================================
\section{Polar H10: recordings and analysis}
% ============================================================

\textbf{Device and protocol:} Polar H10 ECG chest strap, firmware 4.2.0,
JSONL export via Polar Sensor Recording (Android), the author as the only
participant (75~y), 10~September 2026. One spontaneous recording supine
(21.9~min, RR intervals computed on the device), six seated recordings
with a metronome, 5.0--6.0~min each, 130~Hz ECG. Each click marks a change
of respiratory phase; 4, 8, 10 and 12 clicks/min correspond to 2, 4, 5 and
6 breaths/min (0.033 to 0.100~Hz). The two files labelled 4/min are
byte-identical; there is one recording at 4/min. Respiration was not
measured with a separate sensor.

\textbf{Analysis} \KM: R peaks with NeuroKit2 \LM\,[Mak21], artefact
criterion $|\text{RR}/\text{median}-1|>0.30$; flagged intervals are
removed and the remaining beats stay at their true time stamps. Cubic
interpolation to 4~Hz, linear detrending, Hann periodogram over the whole
recording ($\Delta f=1/T\approx0.003$~Hz), band powers by trapezoidal
integration over LF $[0.04,\,0.15)$ and HF $[0.15,\,0.40)$~Hz.

\textbf{Beat detection and time axis.} The simple energy detector in
\texttt{polar\_h10\_atemfrequenz\_scan.py} (fixed threshold at the
90th percentile) found 216--285 beats per recording, the validated
detector 396--443. In every paced recording 35--47\,\% of the beats were
therefore missing. When the remaining intervals are joined end to end,
the time axis shrinks to 49--61\,\% of the recording duration, and all
frequencies shift upwards by a factor of 1.6--2.0. The table is based on
the validated detection with the time axis preserved.

\begin{center}
\small
\begin{tabular}{@{}lrrrrrrrl@{}}
\toprule
Condition & $f_R$~Hz & $T$~min & Beats & RMSSD~ms & LF peak & HF peak & LF/HF & nearest $p/q$ \\
\midrule
Spontaneous, supine & --- & 21.9 & 1517~RR & 19.0 & 0.074 & 0.225 & 1.65 & $33/20$ (0.09\,\%) \\
2/min, rec.\,1 & 0.033 & 5.4 & 431 & 34.2 & 0.067 & 0.200 & 7.44 & --- \\
2/min, rec.\,2 & 0.033 & 5.1 & 411 & 49.6 & 0.067 & 0.167 & 4.24 & --- \\
4/min & 0.067 & 5.3 & 396 & 67.0 & 0.067 & 0.358 & 4.61 & --- \\
5/min, rec.\,1 & 0.083 & 5.6 & 431 & 64.6 & 0.083 & 0.167 & 3.28 & $105/32$ (0.006\,\%) \\
5/min, rec.\,2 & 0.083 & 5.1 & 400 & 68.5 & 0.083 & 0.167 & 2.47 & $121/49$ (0.05\,\%) \\
6/min & 0.100 & 6.0 & 443 & 60.9 & 0.100 & 0.201 & 4.20 & --- \\
\bottomrule
\end{tabular}
\end{center}

\noindent{\small Frequencies in Hz. ---: LF/HF above 4 (outside the grid)
or more than 1.33\,\% from the nearest element. Artefacts: 0--1 each.}

\textbf{Finding 1 --- peak at the breathing frequency} \KM: At 4, 5 and
6/min the dominant peak lies within one frequency bin of the metronome
frequency. Respiration therefore lies in the LF band: the LF peak
\emph{is} the respiratory peak. Where the second harmonic lies above
0.15~Hz (5 and 6/min), it is the HF peak.

\textbf{Finding 2 --- resonance breathing} \KM: RMSSD is 61--69~ms at
4--6/min and 34--50~ms at 2/min. This matches the known amplification of
HRV by slow breathing near the baroreflex resonance \LM\,[Leh14, Sev22].
The spontaneous recording was made supine and is not directly comparable
with the seated recordings.

\textbf{Finding 3 --- 2/min} \SM: The dominant peak lies at 0.067~Hz,
twice the intended breathing frequency, with a smaller component at
0.033~Hz. Without a separate respiration signal it cannot be decided
whether this is the second harmonic of very slow breathing or breathing
at 4/min because of the click convention.

\textbf{Finding 4 --- grid hits at chance level} \KM: The two recordings
at 5/min lie close to two \emph{different} grid elements ($105/32$ and
$121/49$), although the condition is the same; their LF/HF differs by
28\,\%. In 17 sliding 300~s windows of the spontaneous recording, LF/HF
varies between 0.89 and 7.43 (CV 78\,\%); of the 13 windows inside the
grid range $[1,4]$, 6 lie within 0.06\,\% of an element, against a null
rate of 39\,\% (binomial test $p=0.78$). The spread of the estimate is
more than three orders of magnitude larger than the distances to the grid.

\textbf{Spontaneous recording} \KM: mean heart rate 69~/min, RMSSD
19.0~ms, SDNN 48.9~ms. A single missed beat (1741~ms, twice the median)
raises RMSSD to 36.3~ms without artefact handling. Across the 300~s
windows the LF peak varies between 0.040 and 0.085~Hz (CV 22\,\%), far
above 1.33\,\%. LF/HF is 1.65 with one window over the whole recording and
2.21 with 300~s segments; a single estimator setting changes the value by
a third.

\textbf{Methodological lessons:}
\begin{enumerate}
\item Check beat detection against the expected beat count
  (heart rate $\times$ duration) and a validated detector.
\item Preserve the time axis: remove rejected intervals, keep the
  remaining beats at their time stamps, never join them end to end.
\item A single artefact can dominate RMSSD; artefact handling is
  advisable for recordings of any duration.
\item Below 9 breaths/min respiration lies in the LF band; LF/HF then has
  no interpretation as sympathovagal balance.
\item $n=1$ and recordings of 5--6~min allow no statistical conclusions;
  all Polar findings are exploratory.
\end{enumerate}

% ============================================================
\section{Structural reason and consequence}
% ============================================================

\begin{quote}
\textbf{Finding.} \BM\ Under spontaneous breathing the HF peak is the
respiratory frequency (variable, 0.15--0.40~Hz). LF/HF is the ratio of
two independent physiological control loops (baroreflex~/ respiration)
and does not converge to a fixed ratio. With paced breathing the
metronome sets the frequency; peak and harmonics then stand in trivial
integer ratios. In both cases LF/HF is not a test object of the Galois
hypothesis.
\end{quote}

% ============================================================
\section{Next step: prospective test}\label{sec:360next}
% ============================================================

The hypothesis makes a statement about frequency coupling. This can be
tested when the breathing frequency is set deliberately relative to the
individual baroreflex resonance $f_\text{res}$ \SM:

\begin{itemize}
\item Resonance assessment: paced breathing at 4.5--7/min in steps of
  0.5/min, 2~min each \LM\,[Leh14].
\item Four blocks of 21~min (resolution condition at the lowest breathing
  frequency of about 0.062~Hz) at $\rho=f_R/f_\text{res}\in\{2/3,\,3/2\}$
  (grid) and $\{0.618,\,1.618\}$ (golden mean), matched in pairs by
  breathing rate, in random order.
\item Primary outcome: $n{:}m$ phase-synchronisation index
  \LM\,[Tas98] between respiratory phase (respiratory belt) and the HRV
  component at $f_\text{res}$; contrast grid minus golden mean, paired
  two-sided $t$-test.
\item Sample size: $d_z=0.6$, power 0.80 $\Rightarrow n=24$; with
  20\,\% drop-out, 30 participants.
\item Falsification: if coupling is not stronger at grid ratios than at
  the golden mean, the complexity measure $2/(p+q)$ has no basis in HRV.
\end{itemize}

A positive result supports the premise but does not yet distinguish the
prime set and the Farey height of FFGFT from general mode-locking of
coupled oscillators; secondary outcomes serve this purpose (measured
frequency ratio against the grid up to denominator 11, peak stability
against 1.33\,\%).

% ============================================================
\section{Summary}
% ============================================================

\begin{center}
\begin{tabular}{@{}lp{9cm}l@{}}
\toprule
Result & Statement & Status \\
\midrule
Module & \texttt{galois\_hrv\_preprocess.py} works & \KM \\
Band limits & compatible with the 5-limit, not determined by it & \BM \\
Null model & close grid hits of LF/HF are no evidence & \KM \\
BIDMC-01 & negative: too short, intensive-care patient & \LM \\
Fantasia & negative: spontaneous breathing, LF CV 11--33\,\%, 4/10 at chance level & \LM \\
Polar H10 & peak at the breathing frequency, HF = harmonic, hits at chance & \KM \\
Reason & LF/HF is a power ratio; respiration determines the peaks & \BM \\
Next & prospective coupling test, $n=30$, blocks of 21~min & \SM \\
\bottomrule
\end{tabular}
\end{center}

\section*{References}

\begin{description}[leftmargin=2em,labelwidth=1.8em]
\item[{[Pim18]}] Pimentel, M., Johnson, A., Charlton, P., \& Clifton, D.\ (2018).
  BIDMC PPG and Respiration Dataset (v1.0.0). PhysioNet.
  \url{https://doi.org/10.13026/C2208R}
\item[{[Iye96]}] Iyengar, N., Peng, C.\,K., Morin, R., Goldberger, A.\,L., \&
  Lipsitz, L.\,A.\ (1996). Age-related alterations in the fractal scaling of
  cardiac interbeat interval dynamics. \emph{Am.\ J.\ Physiol.\ Regul.\ Integr.\
  Comp.\ Physiol.}\ 271(4), R1078--R1084. Fantasia Database, PhysioNet,
  \url{https://doi.org/10.13026/C2RG61}
\item[{[Leh14]}] Lehrer, P.\,M., \& Gevirtz, R.\ (2014). Heart rate variability
  biofeedback: How and why does it work? \emph{Front.\ Psychol.}\ 5, 756.
\item[{[Sev22]}] Sevoz-Couche, C., \& Laborde, S.\ (2022). Heart rate
  variability and slow-paced breathing: When coherence meets resonance.
  \emph{Neurosci.\ Biobehav.\ Rev.}\ 135, 104576.
\item[{[Mak21]}] Makowski, D.\ et al.\ (2021). NeuroKit2: A Python toolbox for
  neurophysiological signal processing. \emph{Behav.\ Res.\ Methods} 53,
  1689--1696.
\item[{[Tas98]}] Tass, P.\ et al.\ (1998). Detection of $n{:}m$ phase locking
  from noisy data. \emph{Phys.\ Rev.\ Lett.}\ 81, 3291--3294.
\item[{[TF96]}] Task Force ESC/NASPE (1996). Heart rate variability.
  \emph{Eur.\ Heart J.}\ 17, 354--381.
\end{description}

\section*{Scripts}

\texttt{2/python/Dok360\_Skripte/}:
\texttt{galois\_hrv\_preprocess.py},
\texttt{fantasia\_galois\_analyse.py},
\texttt{orbit\_experiment\_synthetic.py},
\texttt{01\_rpeaks\_bidmc01.py},
\texttt{02\_welch\_5limit.py},
\texttt{03\_lombscargle\_dfa.py},
\texttt{pruef\_galois\_nullmodell.py}.
Polar analysis, null model and all numbers of Sections~2 and~5:
\texttt{SWT\_Revision/pruef\_swt\_revision.py} (25/25 assertions).
Fantasia data: f1y01--05, f2o01--05 (.ecg/.hea, CC-BY~4.0).

\section*{Note on the use of AI assistance}

AI assistance (Claude, Anthropic) was used in preparing this document
for wording, LaTeX typesetting and analysis scripts.
The substantive statements and the responsibility for the text
rest with the author.
